\section{A review of \citet{garg2022what}}
% (a) Recap of the main algorithm, its mathematical formulation, and the main thesis of the paper, which you may be questioned on

In this section we review the theory and methods used by \cite{garg2022what}.

\paragraph{Hypothesis.} The hypothesis of the paper is that transformers trained on
specific families of functions, including linear functions, sparse linear functions,
two-layer neural networks, and decision trees, can encode non-trivial learning algorithms
within a single forward pass. Therefore, given in-context examples from a new function
drawn from such a function class, never seen during training, transformers can predict the
function's value at a new query input. ICL is hard to study in LLMs, so the paper invents a
controlled setting to study the phenomenon in transformers only. Concretely for linear
functions, given a set of pairs $\{(x_i, f(x_i))\}_{i=1}^k$ from $f(x)=w^\top x$, a
transformer encodes an algorithm which lets it infer the weights $w$ to predict
$f(x_{k+1})$, resembling least squares. Since every prompt uses a new $w$, memorising a
single function cannot work; the model has to learn how to infer $w$ from the examples.

\paragraph{Formulation.} Let $D_\mathcal{X}$ be a distribution over inputs and
  $D_\mathcal{F}$ a distribution over functions in $\mathcal{F}$. A prompt is a sequence of
  pairs $(x_i, f(x_i))$ followed by a query:
  \begin{equation}
    P =(x_1, f(x_1), \dots, x_k, f(x_k), x_{\text{query}}),
    \label{prompt}
  \end{equation}
with $x_i, x_{\text{query}} \sim D_\mathcal{X}$ i.i.d.\ and $f \sim D_\mathcal{F}$.
A model $M$ \emph{in-context learns} $\mathcal{F}$ up to $\epsilon$ if for some loss
function $\ell$
\begin{equation}
\mathbb{E}_P\!\left[\ell\big(M(P), f(x_{\text{query}})\big)\right] \le \epsilon.
\end{equation}

\paragraph{Training.} A decoder-only GPT-2 style transformer $M_\theta$
(12 layers, 8 heads, 256-dimensional embeddings, 9.5M parameters) is trained
from scratch without any text or pre-trained language model. For
each training step, a function $f$ from the function class and inputs $x_i$ are sampled,
and the model is trained to predict $f(x_{i+1})$ from every prompt prefix
$P^i = (x_1, f(x_1), \dots, x_i, f(x_i), x_{i+1})$, with the objective to minimize
\begin{equation}
\mathbb{E}_P\!\left[\frac{1}{k+1}\sum_{i=0}^{k}
\ell\big(M_\theta(P^i), f(x_{i+1})\big)\right].
\end{equation}

For linear functions, which we study further in this report because of their simplicity
and the particularly interesting ICL behavior, $x_i \sim \mathcal{N}(0, I_d)$,
$f(x)=w^\top x$ with $w \sim \mathcal{N}(0, I_d)$, and $d=20$. Training uses Adam,
500k steps with batch size 64 and a curriculum that gradually increases the input
dimension and prompt length, which significantly speeds up training.

\paragraph{Evaluation.} The transformer trained for ICL is evaluated with respect to the
normalized (by dimension $d$) squared error
\begin{equation}
  \frac{1}{d}(M(P)-f(x_{\text{query}}))^2
  \label{error}
\end{equation}
on prompts $P$ according to \eqref{prompt}, where $k$ is the number of in-context
examples. Its performance in terms of error \eqref{error} is compared to classical
alternatives as a function of the number of in-context examples.

For linear functions, the transformer is compared to minimum-norm least squares, which is
optimal for this noiseless problem, an averaging estimator
$\hat{w}=\frac1k\sum_i x_i y_i$, and 3-nearest neighbours.

\paragraph{Results.}
For linear functions, the trained transformer's error tracks that of least squares for
every number of in-context examples, and is far better than averaging or nearest
neighbours. Performance degrades gracefully under skewed input covariance, a
low-dimensional input subspace, and mismatches between the in-context examples and the
query. It is less robust to rescaling the inputs. The transformer also handles noisy
labels and even reproduces the double descent curve of least squares.

Transformers also learn (i) sparse linear functions, matching Lasso,
(ii) two-layer ReLU networks, matching a network trained by Adam on the context, and
(iii) depth-4 decision trees, outperforming greedy tree learning and XGBoost.

Larger models give lower error and are more robust to distribution shift, and curriculum
learning speeds up training drastically, especially in higher dimensions.